\input{template.tex}

\ifPDFTeX
  \renewcommand{\rmdefault}{cmr}
\fi

\def\SheetNo{3}
\setlength{\emergencystretch}{2em}

\begin{document}
\sheettitle{Множества и отображения.}
\begin{center}\textbf{Краткие решения}\end{center}

\begin{problems}

\item Верно ли, что множество всех функций $f\colon[0,1]\to\mathbb R$ имеет мощность континуума?
\par\smallskip\noindent\textbf{Решение.} Нет; мощность множества таких функций равна $2^{\mathfrak c}$, где $\mathfrak c=|[0,1]|$.
\par Для каждого подмножества $E\subseteq[0,1]$ определим функцию $\chi_E\colon[0,1]\to\mathbb R$: значение $\chi_E(x)$ равно $1$ при $x\in E$ и $0$ при $x\notin E$. Разным подмножествам соответствуют разные функции, поэтому функций не менее $|\mathcal P([0,1])|=2^{\mathfrak c}$.
\par\textit{Другой способ получить нижнюю оценку.} Для каждого $E\subseteq(2,3]$ построим сюръекцию $g_E\colon[0,1]\to E\cup\{2\}$ по формуле
\[
g_E(x)=
\begin{cases}
2+x,&2+x\in E,\\
2,&2+x\notin E.
\end{cases}
\]

\item Бесконечное количество студентов написало контрольную, и каждый получил за неё какой-то балл - число от $0$ до $5$, но никто, кроме практика, не знает баллы за контрольную. Практик им сообщил, что лишь конечное число студентов набрали балл, необходимый для зачёта, и предложил им сыграть в игру: студенты выстраиваются в шеренгу, и каждому на спине напишут его балл за контрольную. Каждый стоящий сзади видит баллы людей, стоящих перед ним, но не видит свой балл и баллы стоящих за ним. Затем с конца шеренги (с человека, который видит всех) практик начинает спрашивать, сколько баллов набрал студент (студент отвечает на ухо практику). Если студент угадает, то получит зачёт, иначе нет, даже если баллов на него хватает. Могут ли студенты сговориться так, чтобы зачёт по итогам игры получило бесконечное количество человек?
\par\smallskip\noindent\textbf{Решение.} Да. Каждый студент называет наименьшую из оценок, которая встречается бесконечно много раз среди видимых ему оценок. Такая оценка существует, поскольку возможных оценок всего шесть. Перед каждым студентом находятся все студенты, кроме него самого и конечного числа стоящих позади. Удаление конечного числа студентов не меняет набора оценок, встречающихся бесконечно часто. Поэтому все называют одно и то же число $c$. Бесконечно много студентов получили именно $c$ и угадывают свою оценку.\par

\item Верно ли, что если для некоторой функции $f$ и множеств $Y_1,Y_2$ выполняется $Y_1\cap Y_2=\varnothing$, то выполняется и $f^{-1}(Y_1)\cap f^{-1}(Y_2)=\varnothing$?
\par\smallskip\noindent\textbf{Решение.} Да. По определению прообраза для любого $x$ из области определения $f$
\[
\begin{aligned}
x\in f^{-1}(Y_1)\cap f^{-1}(Y_2)
&\iff f(x)\in Y_1\ \text{и}\ f(x)\in Y_2\\
&\iff f(x)\in Y_1\cap Y_2=\varnothing.
\end{aligned}
\]
Таких элементов $x$ нет, следовательно, $f^{-1}(Y_1)\cap f^{-1}(Y_2)=\varnothing$.\par

\item Придумайте какое-либо взаимно однозначное соответствие между разбиениями натурального числа на различные и на нечётные слагаемые.
\par\smallskip\noindent\textbf{Решение.} Каждое слагаемое разбиения на различные части единственным образом представим как $2^kq$, где $q$ нечётно. Заменим это слагаемое на $2^k$ экземпляров числа $q$. Получим разбиение того же числа на нечётные слагаемые.
\par Обратно, для каждого нечётного $q$ обозначим через $m_q$ число его вхождений в разбиение и запишем $m_q=\sum_{k\in K_q}2^k$ - двоичное разложение с различными степенями двойки. Заменим $m_q$ экземпляров $q$ на слагаемые $2^kq$, $k\in K_q$. Полученные слагаемые различны: при одинаковом $q$ различаются степени двойки, а равенство $2^kq=2^\ell r$ для нечётных $q,r$ влечёт $k=\ell$ и $q=r$. Оба преобразования сохраняют сумму и взаимно обратны.\par

\item Дана бесконечная последовательность цифр и нечётное число $d$, не делящееся на $5$. Докажите, что можно выбрать несколько цифр подряд, образующих число, делящееся на $d$.
\par\smallskip\noindent\textbf{Решение.} При $d=1$ утверждение очевидно. Пусть $d>1$, а $A_k$ - число, записанное первыми $k$ цифрами последовательности; положим $A_0=0$. Рассмотрим пары остатков $(A_k\bmod d,\,10^k\bmod d)$ для $k=0,1,\ldots,d^2$. Возможных пар не более $d^2$, поэтому найдутся $k<m$ с одинаковыми парами остатков.
\par Из $10^m\equiv10^k\pmod d$ и $\gcd(10,d)=1$ следует $10^{m-k}\equiv1\pmod d$. Если $R$ - число, образованное цифрами с $(k+1)$-й по $m$-ю, то $A_m=10^{m-k}A_k+R$. Поскольку $A_m\equiv A_k\pmod d$, получаем $R\equiv A_m-10^{m-k}A_k\equiv0\pmod d$.\par

\item После того как бесы, коих бесконечно много, проиграли Балде, им стало так стыдно, что они решили заняться физкультурой и записались в спортивные секции. Известно, что в каждую секцию записалось конечное число бесов и среди любого бесконечного набора бесов найдутся двое, занимающиеся в одной и той же секции. Докажите, что бесов-лентяев, которые ходят в конечное число секций, лишь конечное число.
\par\smallskip\noindent\textbf{Решение.} Предположим, что лентяев бесконечно много. Выберем одного из них, $b_1$. Он посещает конечное число секций, в каждой из которых состоит конечное число бесов. Поэтому лентяев, занимающихся хотя бы в одной секции вместе с $b_1$, конечное число. Удалим их и самого $b_1$ из рассматриваемого множества. Останется бесконечно много лентяев.
\par Из оставшихся выберем $b_2$ и повторим процедуру. На каждом шаге удаляется конечное число лентяев, так что выбор можно продолжать неограниченно. Полученная последовательность $b_1,b_2,\ldots$ бесконечна, причём никакие два выбранных беса не занимаются в одной секции: после выбора $b_i$ все его возможные товарищи по секциям удалены. Это противоречит условию задачи. Следовательно, лентяев конечное число.\par

\item Докажите, что следующие множества счётны:
\begin{parts}
  \item множество всех трёхэлементных подмножеств $\mathbb N$;
  \item множество всех конечных подмножеств $\mathbb N$.
\end{parts}
\par\smallskip\noindent\textbf{Решение.} \textbf{а)} Каждому трёхэлементному подмножеству сопоставим тройку его элементов в возрастающем порядке. Это инъекция в $\mathbb N^3$, а $\mathbb N^3$ счётно. Самих подмножеств бесконечно много, поэтому их множество счётно.
\par\textbf{б)} По индукции из счётности $\mathbb N^2$ следует счётность $\mathbb N^k$ при любом $k\ge1$: первые $k-1$ координат кодируются одним натуральным числом, после чего остаётся пара натуральных чисел. Как и в пункте а), множество всех $k$-элементных подмножеств $\mathbb N$ счётно. Все конечные подмножества образуют объединение этих семейств по $k\ge1$ с добавлением пустого множества. Это счётное объединение счётных множеств, следовательно, оно счётно.\par

\item
\begin{parts}
  \item Существует ли континуальная цепочка попарно вложенных друг в друга подмножеств счётного множества?
  \item Может ли иметь континуальную мощность система подмножеств счётного множества, любые два из которых имеют конечное пересечение?
\end{parts}
\par\smallskip\noindent\textbf{Решение.} \textbf{а)} Да. Для каждого $x\in\mathbb R$ положим $S_x=\{q\in\mathbb Q\mid q<x\}$. При $x<y$ имеем $S_x\subsetneq S_y$: между $x$ и $y$ есть рациональное число, принадлежащее $S_y\setminus S_x$. Следовательно, семейство $\{S_x\mid x\in\mathbb R\}$ - континуальная цепочка подмножеств счётного множества $\mathbb Q$.
\par\textbf{б)} Да. Для каждого $x\in(0,1)$ выберем последовательность попарно различных рациональных чисел $q_n(x)\in(0,1)$, сходящуюся к $x$, и положим
\[
T_x=\{0\}\cup\{q_n(x)\mid n\in\mathbb N\}.
\]
Если при $x\ne y$ пересечение $T_x\cap T_y$ бесконечно, то из него можно выбрать последовательность ненулевых общих точек с возрастающими номерами в обеих исходных последовательностях: на каждом шаге их конечные начальные отрезки исключают лишь конечное число точек. Полученная последовательность была бы подпоследовательностью обеих и сходилась бы одновременно к $x$ и $y$, что невозможно.
\par Таким образом, пересечения попарно конечны. Все $T_x$ бесконечны, поэтому при различных $x$ они различны. Семейство имеет мощность континуума, состоит из подмножеств счётного множества $\mathbb Q\cap[0,1]$, причём все они содержат точку $0$.\par

\item Функция $f$ из $\mathbb N$ в $\mathbb N$ сопоставляет числу $n$ наибольший простой делитель числа $n$. Найдите полный прообраз множества чётных чисел.
\par\smallskip\noindent\textbf{Решение.} Ответ: $\{2^k\mid k\ge1\}$. Наибольший простой делитель числа $n\ge2$ чётен тогда и только тогда, когда он равен $2$, то есть когда $n$ является положительной степенью двойки.
\par\textit{Замечание к условию: при $n=1$ наибольший простой делитель не существует. Ответ записан для области определения $n\ge2$.}\par

\item Верно ли, что множество всех непрерывных функций из $\mathbb R$ в $\mathbb R$ континуально?
\par\smallskip\noindent\textbf{Решение.} Да. Зафиксируем перечисление рациональных чисел $q_1,q_2,\ldots$. Каждой непрерывной функции $f\colon\mathbb R\to\mathbb R$ сопоставим последовательность $(f(q_1),f(q_2),\ldots)$. Это соответствие инъективно: если две непрерывные функции совпадают на $\mathbb Q$, то для любого $x\in\mathbb R$ можно выбрать последовательность рациональных чисел $q_{n_j}\to x$ и по непрерывности получить одинаковые значения функций в точке $x$.
\par Следовательно, таких функций не больше, чем последовательностей вещественных чисел, то есть не больше $\mathfrak c^{\aleph_0}=(2^{\aleph_0})^{\aleph_0}=\mathfrak c$. С другой стороны, постоянные функции с различными вещественными значениями различны, а их $\mathfrak c$. Значит, искомая мощность равна $\mathfrak c$.\par

\item Покажите геометрически, что любые два отрезка равномощны и любой интервал равномощен прямой.
\par\smallskip\noindent\textbf{Решение.} Два невырожденных отрезка совмещаются переносом, поворотом и гомотетией; композиция этих преобразований задаёт биекцию.
\par Открытый интервал расположим вдоль открытой полуокружности, сохраняя порядок точек. Проведём касательную к окружности в середине этой дуги. Луч из центра через каждую точку дуги пересекает касательную ровно в одной точке, и каждая точка касательной получается единственным образом. Это задаёт биекцию интервала с прямой.\par

\item В клетки таблицы $3\times6$ записаны числа $1,2,\ldots,9$, каждое число по два раза. Будем называть таблицу хорошей, если в каждой строчке все числа различны и в первых трёх столбцах встречаются все числа от $1$ до $9$. Пусть $N$ - количество хороших таблиц. Найдите степень вхождения $2$ в $N$.
\par\smallskip\noindent\textbf{Решение.} Левый квадрат $3\times3$ заполняется числами от $1$ до $9$ по одному разу, что даёт $9!$ вариантов. Зафиксируем его и обозначим множества чисел в трёх строках через $A,B,C$. В правом квадрате каждое число также должно встретиться ровно один раз; при этом числа из $A$ запрещены в первой строке, из $B$ - во второй, из $C$ - в третьей.
\par Пусть в первую строку правого квадрата попали $k$ чисел из $B$ и $3-k$ чисел из $C$, где $0\le k\le3$. Выбрать их можно $\binom{3}{k}\binom{3}{3-k}$ способами. Оставшиеся $k$ чисел из $C$ обязательно идут во вторую строку. К ним нужно выбрать $3-k$ чисел из $A$, что даёт ещё $\binom{3}{3-k}$ вариантов. Для третьей строки остаются $k$ чисел из $A$ и $3-k$ чисел из $B$, и такое заполнение допустимо.
\par После выбора состава строк числа внутри каждой из них можно переставить $3!$ способами. Поэтому
\[
N=9!\,(3!)^3\sum_{k=0}^{3}\binom{3}{k}^3
=9!\,(3!)^3(1+27+27+1)
=9!\,(3!)^3\cdot56.
\]
В $9!$ двойка входит в степени $\lfloor9/2\rfloor+\lfloor9/4\rfloor+\lfloor9/8\rfloor=7$, в $(3!)^3$ - в степени $3$, а в $56$ - в степени $3$. Искомая степень равна $7+3+3=\boxed{13}$.\par

\item Циклом $(a_1a_2\ldots a_n)$ называется перестановка (биекция на конечном множестве), при которой элемент $a_1$ переходит в $a_2$, в свою очередь, $a_2$ переходит в $a_3$ и так далее до $a_n$, который переходит в $a_1$. Остальные элементы множества переходят сами в себя (остаются на месте). При $n=2$ получается обмен двух элементов, или транспозиция; при $n=1$ получается тождественная перестановка. Покажите, что всякую перестановку можно представить как композицию нескольких непересекающихся циклов (более точно: не пересекаются множества тех элементов, которые циклы не оставляют на месте).
\par\smallskip\noindent\textbf{Решение.} Пусть $f$ - перестановка конечного множества. Выберем элемент $a$ и будем последовательно применять к нему $f$. Поскольку множество конечно, значения $a,f(a),f^2(a),\ldots$ когда-нибудь повторятся. Если $f^s(a)=f^t(a)$ при $s<t$, то, применив обратную перестановку $s$ раз, получим $a=f^{t-s}(a)$. Значит, движение элемента $a$ периодично с самого начала.
\par Пусть $r$ - наименьший положительный период. Тогда элементы $a,f(a),\ldots,f^{r-1}(a)$ различны, а $f^r(a)=a$; они образуют цикл. Если траектории двух элементов встречаются, пусть даже в разные моменты времени, то эти элементы принадлежат одному циклу: из $f^s(a)=f^t(b)$ с помощью обратной перестановки следует, что один элемент получается из другого некоторой степенью $f$. Поэтому разные полученные циклы не пересекаются. Повторяя построение для элементов, ещё не вошедших в циклы, разобьём всё множество на непересекающиеся циклы. Перестановка $f$ совпадает с их композицией.\par

\item По кругу лежат карточки с числами от $1$ до $13$ именно в таком порядке. Изначально все карты лежат числами вверх. Вася может перевернуть карту, лежащую числом вверх, если карта, находящаяся на расстоянии $3$ от данной карты, лежит числом вверх. Например, $12$ можно перевернуть числом вниз, только если $9$ или $2$ лежат числом вверх. Сколько существует способов перевернуть $12$ карт числом вниз? В качестве ответа введите натуральное число.
\par\smallskip\noindent\textbf{Решение.} Будем считать соседними карты, номера которых различаются на $3$ по модулю $13$. Они образуют один цикл, так как $\gcd(3,13)=1$. Во время успешной игры открытые карты должны оставаться связным участком этого цикла. Если бы он распался на две части, то из каждой части пришлось бы оставить хотя бы по одной открытой карте: последнюю карту части перевернуть нельзя, поскольку у неё не будет открытого соседа.
\par Зафиксируем последнюю перевёрнутую карту, например $13$. Оставшейся открытой картой может быть $3$ или $10$ - два варианта. Восстановим игру в обратном порядке. Пока открытый участок содержит от $2$ до $11$ карт, следующую карту можно добавить с любого из двух его концов: на каждом из этих десяти шагов имеется по два варианта. Когда открыто $12$ карт, последняя добавляется однозначно. Значит, при фиксированной последней перевёрнутой карте есть $2\cdot2^{10}=2^{11}$ способов. По симметрии последнюю карту можно выбрать $13$ способами, откуда ответ $13\cdot2^{11}=26624$.
\par\textit{Другой подсчёт.} Перенумеруем карты в порядке $1,4,7,10,13,3,6,9,12,2,5,8,11$. Прибавление $3$ по модулю $13$ обходит все карты, поэтому в этой нумерации условие переворота означает наличие открытой соседней карты. Первую карту можно перевернуть $13$ способами. После этого открытые карты образуют цепочку; чтобы в итоге осталась одна карта, на каждом следующем шаге нужно переворачивать один из двух её концов. Остаётся $11$ таких шагов, каждый с двумя вариантами. Снова получаем $13\cdot2^{11}$.\par

\item Даны $1985$ множеств, каждое из которых состоит из $45$ элементов, причём объединение любых двух множеств содержит ровно $89$ элементов. Сколько элементов содержит объединение всех этих $1985$ множеств?
\par\smallskip\noindent\textbf{Решение.} Ответ: $44\cdot1985+1=87341$. Любые два различных множества пересекаются ровно по одному элементу, поскольку $45+45-89=1$.
\par Зафиксируем одно из множеств $A$. Остальные $1984$ множества пересекают его по одному из $45$ элементов, поэтому некоторый элемент $p\in A$ принадлежит хотя бы $45$ другим множествам. Вместе с $A$ получаем $46$ множеств, попарные пересечения которых равны $\{p\}$.
\par Пусть некоторое множество $B$ не содержит $p$. Оно пересекает каждое из этих $46$ множеств, но само содержит лишь $45$ элементов. Значит, некоторый $q\in B$ принадлежит двум из них. Эти два множества тогда имеют общими и $p$, и $q$, что противоречит размеру их пересечения.
\par Следовательно, все $1985$ множеств содержат $p$, а вне $p$ попарно не пересекаются. Их объединение содержит $1+1985\cdot44$ элементов.\par

\item Докажите, что множество алгебраических чисел, то есть чисел, являющихся корнями многочленов с рациональными коэффициентами, счётно.
\par\smallskip\noindent\textbf{Решение.} Множество $\mathbb Q[x]$ счётно: многочлен степени $d$ задаётся набором из $d+1$ рациональных коэффициентов, а счётное объединение счётных множеств счётно.
\par Каждому алгебраическому числу $\alpha$ сопоставим приведённый многочлен $p_\alpha\in\mathbb Q[x]$ наименьшей степени, корнем которого оно является. Такой многочлен единственен: разность двух различных приведённых многочленов одной минимальной степени имела бы меньшую степень и также обращалась бы в нуль при $\alpha$.
\par У каждого ненулевого многочлена конечное число различных корней. Зафиксируем их нумерацию и обозначим через $j(\alpha)$ номер корня $\alpha$ у многочлена $p_\alpha$. Отображение $\alpha\mapsto j(\alpha)p_\alpha$ инъективно: равенство образов влечёт равенство старших коэффициентов $j(\alpha)$, затем равенство приведённых многочленов и, наконец, самих корней с одинаковым номером. Поэтому алгебраических чисел не более чем счётно. Их бесконечно много, поскольку каждое натуральное число алгебраично.\par

\item У конечного множества $A$ выбрано несколько подмножеств, каждое из которых содержит хотя бы $k^2$ элементов ($k\ge2$). Любые два множества пересекаются, причём не больше чем по $k-1$ элементу. Докажите, что элементы множества $A$ можно покрасить в два цвета так, чтобы каждое из выбранных подмножеств содержало элементы обоих цветов.
\par\smallskip\noindent\textbf{Решение.} Если выбранных подмножеств нет, утверждение очевидно. Иначе выберем одно из них, $S$, и покрасим его элементы в чёрный цвет, остальные - в белый. Каждое другое выбранное множество имеет оба цвета: оно пересекает $S$, но не содержится в $S$, поскольку размер пересечения не превосходит $k-1$.
\par Перекраске элемента $a\in S$ в белый цвет могут мешать только множества $B$, для которых $B\cap S=\{a\}$. Если для какого-то $a$ такого множества нет, эта перекраска сразу даёт требуемый результат.
\par Пусть для каждого $a\in S$ препятствующее множество существует. Выберем различные $a_1,\ldots,a_k\in S$ и соответствующие множества $B_1,\ldots,B_k$ с $B_i\cap S=\{a_i\}$. Покрасим в чёрный цвет множество
\[
R=(S\cup B_1\cup\cdots\cup B_k)\setminus\{a_1,\ldots,a_k\},
\]
а все остальные элементы - в белый. Множества $S,B_1,\ldots,B_k$ имеют оба цвета: из $S$ удалено $k<k^2$ элементов, а из каждого $B_i$ - только $a_i$.
\par Рассмотрим любое другое выбранное множество $C$. В нём есть белый элемент, поскольку
\[
|C\cap R|\le |C\cap S|+\sum_{i=1}^{k}|C\cap B_i|
\le(k+1)(k-1)=k^2-1<|C|.
\]
Если бы в $C$ не было чёрных элементов, то из $C\cap B_i\ne\varnothing$ и $B_i\setminus R=\{a_i\}$ следовало бы $a_i\in C$ для каждого $i$. Тогда $|C\cap S|\ge k$, что противоречит условию. Значит, каждое выбранное множество содержит оба цвета.\par

\item Пусть $X^Y=\{f\mid f\colon Y\to X\}$. Докажите, что существуют биекции между
\[
(X^Y)^Z\ \text{и}\ X^{Y\times Z},
\qquad
(X\times Y)^Z\ \text{и}\ X^Z\times Y^Z.
\]
\par\smallskip\noindent\textbf{Решение.} \textbf{а)} Элемент $h\in(X^Y)^Z$ - это отображение, которое каждому $z\in Z$ сопоставляет функцию $h(z)\colon Y\to X$. Поэтому можно сначала выбрать $z$, получить функцию $h(z)$, а затем вычислить её значение на $y$. Определим
\[
\Phi_1\colon(X^Y)^Z\to X^{Y\times Z},\qquad
\Phi_1(h)(y,z)=h(z)(y).
\]
Обратно, пусть дана функция $F\colon Y\times Z\to X$. При каждом фиксированном $z$ формула $y\mapsto F(y,z)$ задаёт функцию из $Y$ в $X$. Сопоставляя $z$ эту функцию, получаем отображение $\Psi_1(F)\in(X^Y)^Z$:
\[
\Psi_1(F)(z)(y)=F(y,z).
\]
Построения взаимно обратны, поскольку для всех $y\in Y$, $z\in Z$
\[
\Psi_1(\Phi_1(h))(z)(y)=h(z)(y),\qquad
\Phi_1(\Psi_1(F))(y,z)=F(y,z).
\]
Следовательно, $\Phi_1$ - биекция.
\par\textbf{б)} Элемент $f\in(X\times Y)^Z$ - это функция, значением которой при каждом $z\in Z$ является пара из $X\times Y$. Запишем её координаты как $f(z)=(u(z),v(z))$. Тогда $u\colon Z\to X$ и $v\colon Z\to Y$, поэтому определено отображение
\[
\Phi_2\colon(X\times Y)^Z\to X^Z\times Y^Z,\qquad
\Phi_2(f)=(u,v).
\]
Обратно, паре функций $u\colon Z\to X$, $v\colon Z\to Y$ сопоставим функцию $\Psi_2(u,v)\colon Z\to X\times Y$ по правилу
\[
\Psi_2(u,v)(z)=(u(z),v(z)).
\]
Выделение координат пары и составление пары из её координат взаимно обратны: $\Psi_2(\Phi_2(f))=f$ и $\Phi_2(\Psi_2(u,v))=(u,v)$. Следовательно, $\Phi_2$ - биекция.\par

\item Докажите, что квадрат равномощен отрезку.
\par\smallskip\noindent\textbf{Решение.} Достаточно рассмотреть квадрат $[0,1]^2$ и отрезок $[0,1]$. Для каждого числа из $[0,1)$ выберем десятичную запись без бесконечного хвоста девяток, а число $1$ запишем как $0{,}999\ldots$. Паре
\[
x=0{,}x_1x_2\ldots,\qquad y=0{,}y_1y_2\ldots
\]
сопоставим число $F(x,y)=0{,}x_1y_1x_2y_2\ldots$.
\par Если запись результата имеет бесконечный хвост девяток, то такой хвост имеют обе записи координат, откуда $x=y=1$ и $F(x,y)=1$. Во всех остальных случаях запись результата не имеет такого хвоста и потому однозначно определяется числом. Из её цифр на нечётных и чётных местах восстанавливаются записи $x$ и $y$. Следовательно, $F$ - инъекция квадрата в отрезок.
\par В обратную сторону инъекцию задаёт отображение $t\mapsto(t,0)$. По теореме Кантора~- Бернштейна квадрат и отрезок равномощны.\par

\item Докажите, что множество всевозможных бесконечных последовательностей нулей и единиц имеет мощность континуума (то есть равномощно отрезку $[0,1]$).
\par\smallskip\noindent\textbf{Решение.} Последовательности $(a_1,a_2,\ldots)$ сопоставим число
\[
x=\sum_{k=1}^{\infty}\frac{a_k}{2^k}\in[0,1].
\]
Это соответствует последовательному делению отрезка пополам: при $a_k=0$ выбирается левая половина, при $a_k=1$ - правая.
\par Каждое число отрезка имеет такую двоичную запись. Неоднозначность возникает лишь у двоично-рациональных чисел внутри отрезка: одна запись оканчивается нулями, другая - единицами. В каждой такой паре оставим запись, оканчивающуюся нулями. Для концов отрезка оставим последовательности $(0,0,\ldots)$ и $(1,1,\ldots)$ соответственно.
\par Удалённое множество последовательностей счётно: каждая удалённая последовательность задаётся конечным начальным фрагментом, после которого идут только единицы. Оставшиеся последовательности находятся в биекции с $[0,1]$ и имеют мощность $\mathfrak c$. Возвращение счётного множества не меняет эту мощность, поскольку $\mathfrak c+\aleph_0=\mathfrak c$.\par

\item
\begin{parts}
  \item На плоскости изображено бесконечное множество попарно не пересекающихся кругов. Докажите, что это множество счётно.
  \item Назовём восьмёркой фигуру на плоскости, состоящую из двух касающихся внешним образом окружностей. Докажите, что на плоскости можно расположить не более чем счётное множество попарно непересекающихся восьмёрок. Одна восьмёрка может располагаться внутри другой, не должны пересекаться только сами окружности.
\end{parts}
\par\smallskip\noindent\textbf{Решение.} \textbf{а)} Каждый круг содержит непустой открытый квадрат, а значит, точку с рациональными координатами. Зафиксируем нумерацию счётного множества $\mathbb Q^2$ и каждому кругу сопоставим первую по этой нумерации точку, лежащую внутри него. Непересекающимся кругам соответствуют разные точки, поэтому семейство кругов не более чем счётно. По условию оно бесконечно, следовательно, счётно.
\par\textbf{б)} Зафиксируем нумерацию $\mathbb Q^2$. Внутри каждого из двух кругов восьмёрки выберем первую рациональную точку по этой нумерации и сопоставим восьмёрке полученную пару точек, упорядоченную по их номерам. Докажем, что это инъекция.
\par Пусть двум разным восьмёркам соответствует одна пара $(p,q)$. Круги, содержащие $p$, имеют непересекающиеся границы и общую внутреннюю точку, поэтому один строго вложен в другой. Пусть круг второй восьмёрки лежит внутри круга первой. Вторая восьмёрка связна и не пересекает границу этого круга, следовательно, целиком лежит внутри него. Оба ограниченных ею круга также лежат внутри него. Тогда $q$ находится внутри того же круга первой восьмёрки, что и $p$, хотя по выбору $q$ должно лежать внутри другого её круга. Это невозможно, поскольку внутренности двух кругов одной восьмёрки не пересекаются. Значит, соответствие инъективно, а множество пар рациональных точек $\mathbb Q^2\times\mathbb Q^2$ счётно.\par

\item Может ли для некоторой функции $f$ и множеств $X_1,X_2$, таких что $X_1\cap X_2=\varnothing$, быть так, что $f(X_1)=f(X_2)$?
\par\smallskip\noindent\textbf{Решение.} Да. Возьмём $f(x)=x^2$, $X_1=[-2,-1]$ и $X_2=[1,2]$. Тогда $X_1\cap X_2=\varnothing$, но $f(X_1)=f(X_2)=[1,4]$.\par

\item Строка длины $2023$ из $0$ и $1$ называется красной, если в ней есть $7$ единиц подряд. Строка длины $2024$ из $0$ и $1$ называется синей, если в ней есть $8$ одинаковых символов подряд. Каких строк больше, красных или синих, и во сколько раз?
\par\smallskip\noindent\textbf{Решение.} Синей строке $a_1a_2\ldots a_{2024}$ сопоставим строку $b_1b_2\ldots b_{2023}$ по правилу
\[
b_i\equiv a_i+a_{i+1}+1\pmod 2,\qquad 1\le i\le2023.
\]
Иными словами, $b_i=1$, если соседние символы $a_i$ и $a_{i+1}$ совпадают, и $b_i=0$, если они различны. Поэтому восьми одинаковым символам подряд в исходной строке соответствуют семь единиц подряд в новой. Полученная строка красная.
\par Каждая красная строка имеет ровно два прообраза. Действительно, после выбора $a_1\in\{0,1\}$ остальные символы однозначно восстанавливаются: при $b_i=1$ полагаем $a_{i+1}=a_i$, а при $b_i=0$ - $a_{i+1}=1-a_i$. Семь единиц подряд в красной строке обеспечивают восемь одинаковых символов подряд в каждом из двух прообразов, так что оба они синие. Они различны, поскольку начинаются с разных символов. Следовательно, синих строк ровно в $2$ раза больше, чем красных.\par

\item Множество всех целых положительных чисел разбили на два непересекающихся подмножества - $A$ и $B$. Докажите, что для любого целого $n>0$ найдутся такие целые $x>y>n$, что
\[
\{x,y,x+y\}\subseteq A
\]
или
\[
\{x,y,x+y\}\subseteq B.
\]
\par\smallskip\noindent\textbf{Решение.} Зафиксируем $n>0$. Если одно из множеств $A,B$ конечно, то другое содержит все достаточно большие натуральные числа. Возьмём в нём два различных числа $x>y>n$, достаточно больших для того, чтобы и $x+y$ принадлежало этому же множеству.
\par Пусть теперь оба множества бесконечны. Предположим, что требуемой тройки нет. Выберем $y<x<z$ из $A$ так, чтобы $y>n$, $x-y>n$ и $z-x>\max\{n,x-y\}$. Это возможно, поскольку бесконечное подмножество $\mathbb N$ не ограничено сверху.
\par Все три разности $x-y$, $z-x$ и $z-y$ принадлежат $B$. Например, если бы $x-y$ лежала в $A$, то числа $x-y$, $y$, $x$ образовали бы требуемую тройку в $A$; для остальных разностей рассуждение такое же. Но $z-x>x-y>n$ и $(z-x)+(x-y)=z-y$, то есть эти три числа образуют требуемую тройку в $B$. Противоречие.\par

\item  $n \geq 2$ втшников решали $2^{n-1}$ задач. Оказалось, что для каждых двух задач есть ученик, который решил обе эти задачи, и ученик, который решил ровно одну из них. Докажите, что есть задача, которую решили все.
\par\smallskip\noindent\textbf{Решение.} Обозначим множество учеников через $U$, а множество решивших $i$-ю задачу - через $A_i$; положим $B_i=U\setminus A_i$. Множества $A_i$ непусты: иначе для этой и любой другой задачи не было бы ученика, решившего обе. Предположим, что ни одну задачу не решили все. Тогда все $A_i$ и $B_i$ - непустые собственные подмножества $U$.
\par Всего в списке $A_1,B_1,\ldots,A_{2^{n-1}},B_{2^{n-1}}$ имеется $2^n$ множеств, а непустых собственных подмножеств $U$ только $2^n-2$. Поэтому какие-то два множества совпадают. Равенство $A_i=A_j$ при $i\ne j$ невозможно: тогда никто не решил ровно одну из этих задач. Равенство $B_i=B_j$ также невозможно, поскольку влечёт $A_i=A_j$.
\par Остаётся равенство $A_i=B_j$. При $i=j$ множество совпало бы со своим дополнением, что невозможно. При $i\ne j$ получаем $A_i\cap A_j=\varnothing$, то есть никто не решил обе задачи, вопреки условию. Противоречие доказывает, что одну из задач решили все.\par

\end{problems}
\end{document}
